% EUMAS 2026 — Lit Review Expansion Draft
% Four threads for integration into the restructured paper
% Drafted: 2026-04-26 by research agent
% Placeholder \cite{} keys — fill in real bibtex keys before compiling

%% ===================================================================
%% THREAD 1: Prior Art — The DAG Default
%% (For Discussion / Evidence Gap subsection, or new Related Work paragraph)
%% Target: ~200 words
%% ===================================================================

\paragraph{The DAG default.}
The preference for acyclic communication topologies in multi-agent
systems is not an empirical finding---it is a default.
Table~\ref{tab:evidence-gap} documents at least nine major systems
that enforce or favor DAGs ($\beta_1 = 0$) without controlled
evidence that cycles would degrade performance.
Six of the nine---OFA-MAS~\cite{ofamas2026},
G-Designer~\cite{gdesigner2025}, GoAgent~\cite{goagent2026},
ARG-Designer~\cite{argdesigner2026},
AgentConductor~\cite{agentconductor2026}, and
DyTopo~\cite{dytopo2026}---never include cycles in their design
space at all.
MASS~\cite{zhou2025mass} jointly optimizes prompts and topologies but
restricts search to DAGs, then reports that ``prompts matter more than
topology''---a claim that may be an artifact of holding $\kappa = 0$
throughout the search, artificially suppressing the dominant
topological variable our experiments identify.
The Google/MIT scaling study~\cite{googlemit2026scaling} finds 17.2$\times$
error amplification in unstructured communication and a 45\% capability
threshold below which multi-agent coordination degrades performance,
yet frames the solution as ``use fewer agents'' rather than ``choose
better topology.''
Yang et al.~\cite{yang2025topology} explicitly call topology
\emph{the} underexplored dimension of MAS research.
Our confound theorem (Theorem~\ref{thm:confound}) shows why the gap
persists: without density-controlled comparisons, the field cannot
distinguish whether DAGs succeed because acyclicity is beneficial or
because the alternative comparisons are confounded.


%% ===================================================================
%% THREAD 2: Sheaf Upgrade — From beta_1 to H^1(G;F)
%% (For Discussion / Invariant Hierarchy, and Related Work)
%% Target: ~200 words
%% ===================================================================

\paragraph{From graph invariants to sheaf cohomology.}
Our two-invariant framework $(\beta_1, \lambda_2)$ is the
constant-sheaf special case of a richer theory.
Hanks and Riess~\cite{hanks2025cellular} develop cellular sheaves on
graphs, in which each vertex carries a \emph{stalk} (a vector space
representing local agent state) and each edge carries a
\emph{restriction map} encoding how adjacent agents compare information.
For a constant sheaf $\mathcal{F} = \mathbb{R}$ (homogeneous agents),
the first sheaf cohomology group recovers the Betti number exactly:
$\dim H^1(G;\mathcal{F}) = \beta_1$, and the sheaf Laplacian
$L_{\mathcal{F}}$ reduces to the graph Laplacian.
For heterogeneous agents---different models, context windows, or tool
access---the stalks vary across nodes, and the correct invariants
become $(\dim H^1(G;\mathcal{F}), \lambda_2(L_{\mathcal{F}}))$,
capturing topological structure that $\beta_1$ alone misses.

Schmid~\cite{schmid2026prospectus} provides striking independent
validation: a 64-page research prospectus proposing sheaf theory as the
mathematical foundation for multi-agent systems, identifying cellular
sheaves, $H^1$ as coordination obstruction, and sheaf Laplacians as
decentralized controllers.
Schmid's program and ours converge on the same invariants from opposite
directions---decentralized control theory versus compositional
semantics---but Schmid's prospectus lacks the Kleisli formalization,
laxator framework, and experimental grounding that our work provides.
The sheaf upgrade is not merely aesthetic: it predicts that agent
heterogeneity introduces topological degrees of freedom invisible to
$\beta_1$, a testable claim for future density-controlled experiments
with mixed-capability agents.


%% ===================================================================
%% THREAD 3: Information-Theoretic Cost — The Shannon Limit
%% (For Discussion / Invariant Hierarchy, or new subsection)
%% Target: ~200 words
%% ===================================================================

\paragraph{Information-theoretic cost of coordination.}
Thomas and Chen~\cite{thomas2026shannon} establish a lower bound on
multi-agent coordination cost: the minimum number of bits agents must
exchange to achieve consistent global behavior is at least
$\log_2(\dim H^1(G;\mathcal{F}))$.
This upgrades $H^1 = 0$ from a binary possibility test (can agents
coordinate?) to a continuous cost estimator (how expensive is
coordination?), connecting our topological invariants directly to
Shannon information theory.

The bound explains why topology effects scale with problem difficulty
(Experiment~2, Table~\ref{tab:nk-dose}): harder tasks ($K \geq 4$)
demand more coordination, which requires more communication bits,
which means the cohomological floor imposed by $\dim H^1$ becomes
binding. On easy tasks ($K = 0$), the floor is slack and topology
barely matters---exactly what we observe ($\eta^2 = 0.05$ at $K = 0$
versus $\eta^2 = 0.69$ at $K = 6$).

The data processing inequality (DPI) provides a complementary cost
mechanism: Tran and Kiela~\cite{tran2026dpi} show that each agent
boundary in a multi-agent pipeline loses information, with
multi-agent coordination becoming competitive only when input
context is already degraded (noise or masking at $\alpha \geq 0.5$).
Together, these results frame optimal topology selection as
$\beta_1^* = \arg\max[\text{Diversity}(\beta_1) -
\text{Cost}(\beta_1)]$, where the cost combines Shannon capacity
limits and DPI losses.


%% ===================================================================
%% THREAD 4: Certification and Future Directions
%% (For Related Work / Categorical approaches, or Discussion)
%% Target: ~150 words
%% ===================================================================

\paragraph{Categorical certification.}
Capucci and Myers~\cite{capucci2026feedback} formalize a structural
cost of feedback: each cycle in the communication topology requires an
independent convergence certificate, modeled as a categorical
fixed-point condition on traced composition.
This provides the formal cost accounting for our invariant hierarchy:
increasing $\beta_1$ by one gains a unit of transient diversity
(Experiment~4) but incurs one additional convergence obligation.
The cost-benefit framing explains the practical appeal of DAGs---zero
obligations---despite their decision-theoretic suboptimality
\cite{ao2026reliability}: on simple tasks ($K = 0$), the diversity
benefit is negligible and the convergence cost dominates; on complex
tasks ($K \geq 4$), the diversity benefit dominates.
Combined with the companion paper's laxator
formalism~\cite{langer2026games}, this yields the outline of a
complete categorical framework: laxators govern how agents compose
along edges, while the trace governs how composition closes around
cycles.
A full treatment of certification conditions for cyclic multi-agent
composition---connecting Capucci and Myers' convergence guarantees to
our empirical dose-response findings---is deferred to a companion
paper.
